\subsection{Two visible cuts and one private addition}
The remaining addition can either be absorbed into the earlier exposed signals or join the two private products that are not exposed. This distinction gives the last finite obstruction.

\begin{proposition}[Two visible cuts]\label{prop:twovisiblefour}
Under the hypotheses of Proposition~\ref{prop:threevisible}, exactly two of the three deleted additions cannot be visible cuts.
\end{proposition}
\begin{proof}
Apply full-carrier erasure before restriction, and let $q$ be the unique deleted addition that is not visible. The other additions are visible cuts. As before, the exclusion of at most six deletions implies that every original private gate remains and reaches one of their high operands. Except for $q$, a nonproduct operand in this complementary circuit is an old linear clause, possibly with earlier visible high signals inserted. Full quadratic carriers have been erased from its inputs.

First suppose $q$ itself is a cut's high operand. It has minimum degree at least two, so both its operands do also. Neither operand can be old linear or a degree-one visible cut; a degree-two visible cut is a protected carrier and is also excluded. Thus $q=u\vee v$ for two private products, possibly identical. Replace every Boolean occurrence of $q$ by $u\vee v$ and expose $u,v$ at their creation times. The second visible high operand is another private product or $q$ itself. There are now at most three exposed product outputs. Their direct contexts are the original permitted contexts, duplicated when needed, and pair terms expand to pair terms; repeated squares are absorbed into constant-one contexts. A private operand previously equal to $q$ becomes the permitted linear-mode operand with zero old part and exposed signals $u,v$. Thus the function belongs to the three-high or two-high relaxation already excluded. Every product is ancestral to a designated high, so the last product remains designated.

We may therefore assume that both visible highs are private products $g,f$. If both operands of $q$ are old or perturbed linear clauses, their union is another permitted linear-mode operand, giving the two-high relaxation. If its operands are a private product $u$ and a linear clause, expose $u$ as a third signal, with zero direct context and no new output pair terms. The value of $q$ is then a permitted linear-mode operand. The same argument applies when one of two product operands is already designated: include that signal in the linear part and expose the other one if necessary. Repeated identical operands need no separate exposure. These reductions enlarge the Boolean class; they do not assert that the original arithmetic circuit can lose an addition while preserving its extraction.

If $g=f$, expose both product operands of $q$ as two further signals, with zero direct contexts and no new output pair terms. There are at most three designated products, and $q$ is their permitted linear-mode sum. Every product still reaches the designated $g$, so the three-high relaxation applies. Thus only $q=u\vee v$ with two distinct undesignated private products and distinct $g,f$ remains. The four products are exactly $u,v,g,f$. Every product is ancestral to $g$ or $f$, so the last topological product is designated. Numbering the products from one to four, their designated sets are
\[
 \{1,4\},\qquad\{2,4\},\qquad\{3,4\}.
\]
The other two products are the operands of $q$. Allow $q$ immediately after both have been created. Each product operand is an old linear clause OR any subset of preceding designated highs, a previous product alone, or $q$ alone when available. The restricted output is
\[
 gL\vee fM\vee\delta gf,
 \qquad\delta\in\{0,1\},
\]
with the same allowed affine contexts as in Proposition~\ref{prop:threevisible}.

The three exceptional templates are again finite over nine inputs, after identifying all isolates with one variable. Select the stated input modes, all old variable memberships, all earlier designated-high memberships, both admissible affine contexts, and the pair term independently. Empty old clauses and zero contexts are allowed. No restricted gate is required to be nonzero, distinct or essential. Exact Tseitin equivalences impose truth on the four edge assignments and falsity on the sixteen maximal independent transversals; monotonicity makes these tests complete.

The three CNFs have $(3028,10563)$, $(2986,10403)$ and $(2986,10409)$ variables and clauses. Their complete DRUP proofs have $47\,892$, $54\,127$ and $93\,835$ lines. Independent regeneration reproduced the CNFs exactly, and independent RUP-only checking verified the refutations, with zero RAT steps, in $0.97$, $1.07$ and $2.08$ seconds. The complete data and hashes accompany the visible-cut obstruction. Hence all three exceptional classes are impossible, completing the proof.
\end{proof}

\begin{theorem}[The optimal seven-deletion case]\label{thm:sevendeletion}
Suppose exactly four quadratic witnesses are missing, $d=4$, and the extracted quadratic circuit is optimal. Then seven deletions are impossible. Consequently a circuit below the star bound with four missing witnesses must satisfy
\[
\begin{array}{c|l}
d&(D,e)\\\hline
2&(5,0)\\
3&(5,0),(5,1),(6,0)\\
4&(6,1).
\end{array}
\]
In the last row exactly four products and two additions are deleted.
\end{theorem}
\begin{proof}
Proposition~\ref{prop:fiveproductseven} leaves only four products and three additions. At least one addition is a visible cut; otherwise the selected degree-two output would equal the full output and no quadratic witness would be missing. The one-, two- and three-visible-cut cases are excluded by Propositions~\ref{prop:onevisible}, \ref{prop:twovisiblefour} and~\ref{prop:threevisible}. The displayed alternatives follow from the earlier four-witness reduction and the six-deletion exclusion.
\end{proof}
